\documentclass[conference]{IEEEtran}
\usepackage[utf8]{inputenc}
\usepackage[T1]{fontenc}
\usepackage[spanish]{babel}
\usepackage{cite}
\usepackage{booktabs}
\title{Evolución funcional de una Plataforma de Instrumentación Didáctica: diagnóstico y requerimientos de PID 2.0}
\author{\IEEEauthorblockN{Carlos Alberto Becerra Grajales y Miriam Crisanto Zetina}\IEEEauthorblockA{Doctorado en Sistemas Computacionales\\Universidad del Sur}}
\begin{document}
\maketitle
\begin{abstract}
Se presenta la integración preliminar de las sesiones 1 y 2 de una investigación sobre la evolución de la Plataforma de Instrumentación Didáctica (PID). La inspección estática del sistema original documentó carga de instrumentaciones XLSX, sustitución de archivos, fechas de entrega y reportes. Se diseñaron una entrevista semiestructurada y una guía de observación; esta última fue aplicada parcialmente a código y esquema SQL. Se formularon 12 requerimientos funcionales y siete no funcionales con trazabilidad documental. Un conjunto de datos sintéticos define escenarios futuros de aprobación y rechazo de cambios. No se realizaron entrevistas ni pruebas funcionales, y los requerimientos continúan sujetos a validación.
\end{abstract}
\begin{IEEEkeywords}instrumentación didáctica, ingeniería de requerimientos, trazabilidad, control de cambios\end{IEEEkeywords}
\section{Introducción y problema de investigación}
La instrumentación didáctica organiza contenidos, actividades, evaluación y tiempos en la educación superior tecnológica. El PID original administraba entregas documentales de instrumentaciones y reportes. La última versión recibía instrumentaciones XLSX, aunque algunos borradores de la sesión 1 describían incorrectamente PDF como formato de entrada. El sistema también contempla reportes de salida PDF.

El problema de investigación consiste en determinar cómo evolucionar del registro documental al seguimiento estructurado por unidad temática y actividad, con fechas programadas y reales, solicitudes de cambio autorizadas y conservación de versiones.

\subsection{Pregunta y objetivo}
¿Cómo diseñar un modelo de mejora funcional del PID que integre seguimiento programático y gestión autorizada de modificaciones a la planeación docente, fortaleciendo control y trazabilidad? El objetivo general es analizar y diseñar un modelo de mejora funcional con seguimiento por unidad y actividad, calendarios y modificaciones autorizadas.

\subsection{Alcance y continuidad de la sesión 1}
El objeto de estudio comprende planeación, seguimiento y cambios; excluye control escolar, calificaciones individuales y finanzas. La revisión crítica de antecedentes, análisis comparativo de ciclo de vida, factibilidad y referencias completas de P1 están pendientes de integración fiel desde su versión definitiva; este apartado es un puente editorial, no reemplaza P1.

\section{Metodología de la segunda sesión}
El levantamiento se orientó por los principios de ingeniería de requerimientos y trazabilidad \cite{nuseibeh2000,wiegers2013}. Se inspeccionó estáticamente el código y esquema SQL depurados de PID-Original. Se diseñaron I-01 (entrevista semiestructurada de 16 preguntas) e I-02 (guía de observación de 14 ítems). I-02 se aplicó parcialmente a archivos remotos del repositorio; I-01 no se ha aplicado.

Se distinguieron evidencias de código (EV-COD), inspección (EV-OBS), aclaraciones históricas (CTX), propuestas (PROP) y simulaciones (SIM). Las afirmaciones negativas se limitan a los archivos revisados. Se diseñó un dataset completamente ficticio, SIM-001, para preparar escenarios de prueba, no para presentar resultados empíricos.

La calidad preliminar de los requerimientos se valoró mediante identificación, fuente, prioridad, aceptación verificable, consistencia, factibilidad y trazabilidad. La validación con actores académicos permanece pendiente.

\section{Resultados preliminares}
\subsection{Inspección del sistema original}
El esquema depurado contiene 16 tablas, incluyendo instrumentaciones, archivos, usuarios, materias, periodos y permisos. En el controlador se identificó carga XLSX, registro de fecha, contador de modificaciones y cuatro operaciones de entrega de reportes. En las funciones revisadas no se acreditó edición estructurada de actividades ni un flujo íntegro de aprobación y versionado inmutable; esto no demuestra su ausencia en todo el sistema.

Se identificaron riesgos de seguridad por comparación directa de credenciales, generación de sesión con MD5 y uso de identificador de sesión por URL. La revisión estática no constituye prueba de explotación.

\subsection{Especificación y trazabilidad}
Se formularon 12 requerimientos funcionales (RF-01 a RF-12), siete no funcionales (RNF-01 a RNF-07) y cinco reglas candidatas. La matriz documental vincula los 19 requerimientos con necesidad, objetivo provisional, fuente y prueba prevista. La cobertura documental 19/19 no implica validación empírica.

\subsection{Diccionario y escenarios sintéticos}
Se inventariaron 16 tablas originales y propusieron 13 entidades conceptuales. SIM-001 describe una solicitud aprobada que genera versión 2 y una rechazada que mantiene la versión vigente. Estas decisiones son datos simulados predefinidos; no son salidas de software ejecutado.

\section{Discusión preliminar y limitaciones}
El aporte potencial no consiste únicamente en incorporar interfaces, sino en especificar un modelo reproducible de seguimiento programático, control autorizado de cambios y trazabilidad. La inspección del código proporciona evidencia técnica del antecedente, pero no prueba el beneficio de PID 2.0.

No se han ejecutado entrevistas, pruebas de aceptación ni observación funcional. Permanecen pendientes la autoridad competente para aprobar cambios, la fórmula de avance, la identidad de actividades entre versiones y la evaluación de seguridad. En la sesión 3 se prevé resolver el modelo de datos y flujos; en la sesión 4, implementar y evaluar una PoC delimitada con datos sintéticos. No se formulan conclusiones finales.

\bibliographystyle{IEEEtran}
\bibliography{referencias}
\end{document}
